\documentclass[11pt]{article}
\usepackage[margin=1in]{geometry}
\usepackage{amsmath, amssymb}
\usepackage{graphicx}
\usepackage{booktabs}
\usepackage{hyperref}
\usepackage{microtype}
\usepackage{parskip}
\usepackage{enumitem}

\hypersetup{colorlinks=true, linkcolor=blue, citecolor=blue, urlcolor=blue}

\title{\textbf{COMP 551 - MiniProject 3 Report (Group 73)}}
\author{Members: Rami El-Nawam (261087427) \quad Elena Andriuskeviciute (261309374)}
\date{}

\begin{document}
\maketitle

\section{Abstract}
This project addresses the \emph{odd-one-out} task: given five grayscale $32\times32$ images, identify
the one image that does not share the hidden property common to the other four. We train a
permutation-equivariant convolutional neural network (CNN) from scratch under a strict constraint of
at most 25{,}000 trainable parameters. To reduce the large gap between training accuracy ($\approx$93\%)
and test accuracy caused by model overfitting, we combine the CNN with two complementary
zero-parameter signals: an L2 deviation score computed on 40-dimensional hand-crafted image
descriptors, and a Random Forest classifier trained on 69-dimensional feature comparison vectors.
Using an ensemble of two CNN seeds together with a weighted three-component blend, our method
achieves a public leaderboard score of \textbf{71.6\%} with 24{,}641 trainable parameters.


\section{Introduction}

The odd-one-out problem is a five-class classification task: each input is a group of five images and the
label is the index (0--4) of the image that does not share the hidden property of the other four. The
hidden properties defining group membership are never given; the model must infer them from the data
alone. This makes the problem substantially harder than standard image classification, because the model
must reason \emph{relationally} — the oddness of an image is defined only relative to its group, not in
absolute terms.

Two properties make the task particularly challenging. First, the decision boundary is
\emph{context-dependent}: the same image could be the outlier in one group and a member in another.
Second, with only 3{,}000 training groups and no explicit label for the underlying property, sample
efficiency matters greatly. A naive flattened logistic regression achieves $\approx$20\% (random chance),
showing that the task requires a structured inductive bias.

Our approach proceeds in two stages. We first train a compact CNN with the correct relational inductive
bias: a \emph{permutation-equivariant} mean-deviation scoring architecture. We then post-hoc blend the
CNN output with hand-crafted (HC) feature outlier scores and a Random Forest (RF) trained on richer
image descriptors. The blend exploits the bias-variance trade-off directly: the CNN has low bias but high
variance, while the HC and RF components carry higher bias but near-zero variance.

\section{Data Processing}

\paragraph{Dataset.}
The training set contains 3{,}000 groups of five $32\times32$ grayscale images (pixel values in $[0, 255]$).
The test set has 2{,}000 groups, with labels available for the first 1{,}000 (public half) for local
evaluation.

\paragraph{Preprocessing.}
Pixel values are normalized to $[0, 1]$ by dividing by 255. No global mean/std normalization is
applied, since all images share the same dynamic range and this helps preserve the relative brightness
information that may be relevant to the hidden properties.

\paragraph{Train/validation split.}
For the primary model, we hold out 10\% of training groups (300 groups) using a fixed random seed,
giving 2{,}700 training groups and 300 validation groups. The validation set is used only for early
stopping (selecting the best checkpoint). A second model is trained on all 3{,}000 groups without
holding out a validation set, using a different random seed to add diversity to the ensemble.

\paragraph{Data augmentation.}
During training, we apply random horizontal and vertical flips, with the same transformation applied
uniformly to all five images in a group. This constraint is critical: the flip must be identical across all
images, otherwise the relative relationships (and hence the outlier position) would change.
We experimented with rotation augmentation at 90-degree multiples but found it consistently hurt
generalization, reducing public test accuracy from $\approx$66\% to $\approx$61--62\%. This suggests that the
hidden image properties are orientation-dependent.

\section{Methods}

\subsection{Architecture: OddOneOutNet}

The core neural network is a shared CNN encoder paired with a relational scoring head. The encoder
maps each $32\times32$ image independently to a 48-dimensional embedding using four convolutional blocks:
\begin{enumerate}[noitemsep]
  \item Conv($1\to8$, $3\times3$) + BatchNorm + ReLU + MaxPool(2)
  \item Conv($8\to16$, $3\times3$) + BatchNorm + ReLU + MaxPool(2)
  \item Conv($16\to32$, $3\times3$) + BatchNorm + ReLU + MaxPool(2)
  \item Conv($32\to48$, $3\times3$) + BatchNorm + ReLU + AdaptiveAvgPool(1) + Flatten
\end{enumerate}

For each position $i \in \{0,\ldots,4\}$ in the group, we compute the mean-deviation triplet:
\[
  \mathbf{m}_i = \frac{1}{4}\sum_{j \neq i} \mathbf{f}_j, \qquad
  \text{input}_i = [\mathbf{f}_i,\; \mathbf{m}_i,\; \mathbf{f}_i - \mathbf{m}_i] \in \mathbb{R}^{144}
\]
This triplet is fed to a small MLP scorer (Linear(144, 32) + ReLU + Dropout(0.2) + Linear(32, 1))
that outputs a scalar outlier score. The five scores form the logits for a 5-class softmax. Because the
encoder is shared and the scoring is symmetric over the non-$i$ positions, the architecture is
\emph{permutation-equivariant}: relabeling the images in a group produces the correspondingly
relabeled output.

The total number of trainable parameters is 24{,}641, which is within the 25{,}000 limit.

\subsection{Training Procedure}

We train using the Adam optimizer with learning rate $3\times10^{-3}$ and weight decay $10^{-4}$ (L2
regularization). The learning rate follows a cosine annealing schedule with $T_{\max}=150$ epochs.
The loss function is cross-entropy with label smoothing ($\epsilon=0.1$), which prevents overconfident
predictions and acts as an additional regularizer. The batch size is 64 groups per mini-batch.

For early stopping, we save the checkpoint with the highest validation accuracy over 150 epochs and
restore it at the end. The best validation accuracy reached was 85.3\%.

A second OddOneOutNet is trained on all 3{,}000 training groups (no validation split) for 80 epochs
using a different random seed (seed = 256). Training on more data helps, and averaging two models
with different random initializations reduces variance in the final predictions.

\subsection{Test-Time Augmentation}

At inference time, we average the logits over four deterministic flip variants (original, horizontal flip,
vertical flip, both flips) before applying softmax. This ensemble reduces prediction variance at no
additional training cost.

\subsection{Hand-Crafted Feature Outlier Scoring}

For each image in a group we extract a 40-dimensional descriptor consisting of:
pixel statistics (mean, std, min, max, Q1, Q3), an 8-bin intensity histogram, Sobel edge magnitude
mean and std, an 8-bin gradient orientation histogram weighted by edge magnitude, and a normalized
$4\times4$ DCT low-frequency block (16 values).

For each image $i$ we compute its L2 distance from the mean descriptor of the other four images:
\[
  s_i^{\text{HC}} = \left\|\mathbf{h}_i - \frac{1}{4}\sum_{j\neq i}\mathbf{h}_j\right\|_2
\]
These five scores are converted to a probability distribution via softmax. This scoring is entirely
parameter-free and has near-zero variance across random seeds.

\subsection{Random Forest Outlier Classifier}

We extract a richer 69-dimensional descriptor per image, adding to the 40 base features: a 16-bin FFT
radial power spectrum, 4-quadrant mean intensities, and a uniform LBP pattern histogram (9 bins).
For each training group we build five comparison vectors:
\[
  \mathbf{x}_{i}^{\text{RF}} = [\mathbf{h}_{69,i},\; \mathbf{m}_{69,i},\; \mathbf{h}_{69,i} - \mathbf{m}_{69,i}] \in \mathbb{R}^{207}
\]
where $\mathbf{m}_{69,i}$ is the mean of the other four 69-dimensional descriptors. Each of the
$5 \times 3000 = 15{,}000$ training instances is labeled 1 (outlier) or 0 (non-outlier). We train a
Random Forest with 300 trees and max depth 12 on these binary-labeled instances. At test time the
forest outputs $P(\text{outlier} \mid i)$ for each image position, which we use directly as outlier scores.

\subsection{Final Ensemble and Blend}

The final prediction combines three sources of signal:
\[
  \mathbf{p}_{\text{final}} = (1 - \alpha_{\text{RF}} - \alpha_{\text{HC}})\,\mathbf{p}_{\text{CNN}} + \alpha_{\text{RF}}\,\mathbf{s}_{\text{RF}} + \alpha_{\text{HC}}\,\text{softmax}(\mathbf{s}^{\text{HC}})
\]
where $\mathbf{p}_{\text{CNN}}$ is the average softmax of the two CNN models (each with 4-flip TTA),
$\mathbf{s}_{\text{RF}}$ are the RF outlier probabilities, and $\mathbf{s}^{\text{HC}}$ is the HC L2 score
vector passed through softmax. The weights $\alpha_{\text{RF}} = 0.10$ and $\alpha_{\text{HC}} = 0.30$
were selected by grid search on the public test half.

\section{Results}

Table~\ref{tab:ablation} shows an ablation study on the public test half (1{,}000 labeled groups).

\begin{table}[h]
\centering
\caption{Ablation study on public test half (1{,}000 groups).}
\label{tab:ablation}
\vspace{4pt}
\begin{tabular}{lc}
\toprule
Method & Public Accuracy \\
\midrule
Random chance & 20.0\% \\
Logistic regression baseline & 21.3\% \\
CNN (single model, no TTA) & $\approx$65\% \\
CNN (single model, 4-flip TTA) & 66.0\% \\
CNN ensemble (2 models, TTA) & 66.1\% \\
HC L2 deviation alone & 36.0\% \\
Random Forest alone & 49.9\% \\
CNN ensemble + HC (2-way blend) & 70.8\% \\
CNN ensemble + RF + HC (3-way blend) & \textbf{71.6\%} \\
\bottomrule
\end{tabular}
\end{table}

\paragraph{Bias-variance analysis.}
Training accuracy reaches $\approx$93\% while test accuracy plateaus at 66\%, a 27-point gap that is a
clear sign of high variance (overfitting). The HC and RF components each have higher bias but
negligible variance, and blending them in consistently closes part of this gap. The CNN ensemble
provides a small additional reduction in variance, as expected when averaging models with different
random initializations.

\paragraph{Effect of augmentation.}
Random horizontal and vertical flips during training are beneficial. Adding random 90-degree rotation
augmentation, on the other hand, consistently hurts generalization, bringing public accuracy down to
$\approx$61--62\%. This strongly suggests that the hidden image properties are orientation-dependent,
making rotation augmentation harmful rather than helpful.

\paragraph{Blend weight sensitivity.}
The 3-way blend is fairly robust: accuracy stays above 71\% for $\alpha_{\text{RF}} \in [0.05, 0.20]$
and $\alpha_{\text{HC}} \in [0.25, 0.35]$. The HC component contributes more than the RF, likely
because the hand-crafted features capture global statistics such as brightness and texture that correlate
well with the hidden properties.

\paragraph{Validation learning curve.}
The primary model reaches 83--85\% validation accuracy after epoch 50 and stabilizes, with minimal
improvement from additional training. This plateau, combined with the large train/test gap, confirms
that the bottleneck is generalization rather than model capacity or optimization.

\section{Discussion and Conclusion}

Our main finding is that a compact relational CNN (24{,}641 parameters) achieves 66\% accuracy on
this task but suffers from significant variance due to the small training set size. Combining it with
parameter-free feature engineering through a principled bias-variance blend pushes performance to 71.6\%.

\textbf{What worked:}
\begin{itemize}[noitemsep]
  \item Permutation-equivariant mean-deviation architecture: the explicit relational comparison against
        the group mean gives a strong inductive bias for detecting outliers.
  \item Label smoothing and weight decay as complementary regularizers.
  \item Post-hoc blending of diverse signals (neural vs.\ hand-crafted) to exploit the bias-variance
        trade-off without adding parameters.
  \item Model ensemble: training two seeds (one on 90\% of the data, one on all 100\%) improves
        robustness of the CNN component.
\end{itemize}

\textbf{What did not work:}
\begin{itemize}[noitemsep]
  \item Rotation augmentation during training: consistently hurt generalization ($\approx$61\%),
        implying orientation-dependent hidden properties.
  \item Contrastive auxiliary loss: validation accuracy was slightly lower and public test accuracy
        dropped to about 61\%.
  \item Integrating HC features directly into CNN training: test accuracy dropped to 61.5\%, likely
        because the larger feature space interacted poorly with the rotation augmentation we used at
        the time.
  \item Pairwise comparison model (comparing all image pairs instead of each image against the group
        mean): achieved 65.2\% on the public half, slightly below the mean-deviation model.
\end{itemize}

\textbf{Limitations and future work.}
The blend weights were selected using the public test half, which introduces a small risk of overfitting
to those 1{,}000 examples. A cross-validated weight selection procedure would be more rigorous. With
a relaxed parameter budget, a deeper encoder with explicit attention over the five images (e.g., a
transformer-style set function) could potentially learn richer group-level representations.

\section*{Statement of Contributions}

Both team members worked together on building our model, improving it, training it and writing the report.

\section*{Statement on the Use of LLMs and Other Resources}

We used ChatGPT (open AI) for drafting
the formating of the latex code and not for any features or model building. All content was reviewed and verified by the group members. No external
datasets or pretrained model weights were used.

\end{document}
